Extra-classical receptive field (eCRF) effects, such as surround suppression, demonstrate that local processing within the primary visual cortex (V1) is insufficient for global visual integration. While hierarchical models incorporating feedback have been proposed, simple feedback architectures often fail to replicate the robust suppression observed physiologically. This study proposes that hierarchical competitive dynamics between cortical areas are essential for the emergence of these effects. Using the NEST simulator, we developed a large-scale spiking neural network based on the canonical cortical microcircuit, integrating layered V1 and V2 modules. We systematically evaluated three configurations: (1) V1 lateral connections only, (2) non-specific V2 feedback, and (3) a hierarchical framework with lateral competition between V2 modules. Our simulations reveal that neither V1 lateral connections nor simple feedback loops generate physiological surround suppression. Crucially, robust suppression (~75\%) consistent with experimental data was only achieved through the integration of vertical pathways with lateral competition between V2 modules. These findings suggest that eCRF phenomena are emergent properties of inter-areal competitive dynamics rather than purely local or linear hierarchical processes. Furthermore, we demonstrate how alterations in the excitation/inhibition balance within this framework replicate gamma-band oscillations associated with neuropsychiatric disorders. Our results highlight the importance of hierarchical competition in sensory processing and provide a bridge between circuit dynamics and computational psychiatry.

